% Part: turing-machines 
% Chapter: undecidability
% Section: introduction

\documentclass[../../../include/open-logic-section]{subfiles}

\begin{document}

\olfileid{tur}{und}{int} 
\olsection{تعارف}

ایہ گل شاید صاف لگے
کہ کجھ فنکشن، ایتھوں تک کہ
کجھ حسابی فنکشن وی،
حساب کرن جوگ نہیں۔
فنکشن بہت
ودھ نیں تے اوہناں
دا چلن بہت پیچیدہ
ہو سکدا اے۔ مثال
دے طور اُتے
تابکار ذریاں دے
ٹُٹّن، یا ہور
بے قاعدہ یا اتفاقی
چلن نال تعریف
ہون والے فنکشن
ویکھو۔ فرض کرو
اسی کسے مقرر ویلے
توں اک اک سیکنڈ
دے وقفے گننا
شروع کریے،
تے فنکشن $f(n)$
نوں کائنات دے
اوہناں ذریاں دی
گنتی رکھئیے جیڑے
اُس ابتدائی ویلے
توں بعد $n$-ویں
اک سیکنڈ دے
وقفے وچ ٹُٹدے
نیں۔ ایہ ایسا
فنکشن لگدا اے
جس دا حساب
کرن دی امید
اسی کدی نہیں
رکھ سکدے۔

پر ایہ تصور نہ کر
سکنا کہ ایہو جیہے فنکشن
دا حساب کِویں
کیتا جاوے، اوہنوں
واقعی ناقابلِ حساب
ثابت کرن توں
بالکل وکھ گل اے۔
حقیقت وچ بظاہر
انتہائی پیچیدہ
فنکشن وی تجریدی
معنیٰ وچ حساب
ہو سکدا اے۔
مثلاً فرض کرو
کائنات دا وقت
محدود اے:
کجھ کائناتی
نظریاں مطابق
بہت دور مستقبل
وچ کائنات اک
نکتے وچ سکڑ
جاوے گی۔ فیر
ابتدائی ویلے توں
بعد صرف محدود،
بھاویں ناقابلِ
تصور حد تک
وڈی، تعداد وچ
سیکنڈ نیں جنہاں
لئی $f(n)$
تعریف شدہ اے۔
تے جیہڑا فنکشن
صرف محدود
انپٹاں اُتے
تعریف شدہ ہووے،
اوہ حساب کرن
جوگ اے: اسی
اوہدے آؤٹ پٹ
اک وڈے جدول
وچ لکھ سکدے
آں، یا اک بہت
وڈے ٹیورنگ
مشین دے حالتی
تبدیلی والے
خاکے وچ کوڈ
کر سکدے آں۔

اکثر سانوں اوہناں
فنکشناں دیاں
خاص صورتاں وچ
دلچسپی ہوندی اے
جنہاں دیاں قدراں
ہاں یا ناں والے
سوالاں دے جواب
دیندیاں نیں۔
مثلاً “کی $n$
اولی عدد اے؟”
سوال نال ایہ
فنکشن جڑیا اے:
\[
\fn{isprime}(n) = \begin{cases}
  1 & \text{جے $n$ اولی عدد ہووے}\\
  0 & \text{ہور صورت وچ۔}
  \end{cases}
\]
اسی کہندے آں کہ
ہاں/ناں والے
سوال دا \emph{مؤثر
طریقے نال فیصلہ}
ہو سکدا اے جے
اوہدے نال جڑیا
$1/0$-قدری
فنکشن مؤثر
طریقے نال حساب
ہو سکدا ہووے۔

ریاضیاتی طور اُتے
ایہ ثابت کرن لئی
کہ کجھ فنکشن
مؤثر طریقے نال
حساب نہیں ہو
سکدے، یا کجھ
مسئلیاں دا مؤثر
طریقے نال فیصلہ
نہیں ہو سکدا،
حساب دا اک خاص
نمونہ طے کرنا
ضروری اے۔ فیر
ایہ دکھانا پیندا
اے کہ کجھ فنکشن
اوہ نمونہ حساب
نہیں کر سکدا،
یا کجھ مسئلیاں
دا فیصلہ نہیں
کر سکدا۔ مثلاً
اسی دکھا سکدے
آں کہ کجھ فنکشن
ٹیورنگ مشین نال حساب
نہیں ہو سکدے تے کجھ
مسئلیاں دا فیصلہ
ٹیورنگ مشین نہیں کر سکدی۔
فیر چرچ--ٹیورنگ
دے دعوے دا حوالہ
دے کے ایہ نتیجہ
کڈھ سکدے آں کہ
صرف ٹیورنگ
مشیناں ہی ہر
فنکشن دا حساب
کرن توں قاصر
نہیں، بلکہ کوئی
وی مؤثر طریقۂ
کار ایہ کم نہیں
کر سکدا۔

ایہناں منفی نتیجیاں
نوں ثابت کرن دی
اہم گل ایہ اے
کہ اسی خود ٹیورنگ
مشیناں نوں عدد
دے سکدے آں۔
ایس دا سوکھا
طریقہ اوہناں دی
گنتی وار فہرست
بناؤنا اے: مثلاً
ٹیورنگ مشیناں
تے اوہناں دے
پروگرام لکھن دا
اک مقرر طریقہ
چن کے فیر
اوہناں نوں
باقاعدہ ترتیب
نال فہرست وچ
لکھو۔ جدوں
ایہ ممکن ہونا
صاف ہو جاوے،
تاں سیٹاں دے سائز
دی سادی دلیل نال
ٹیورنگ مشین
توں ناقابلِ
حساب فنکشناں
دا وجود نکلدا
اے: $\Nat$ توں~$\Nat$
فنکشناں دا سیٹ
(بلکہ صرف
$\Nat$ توں
$\{0, 1\}$
فنکشناں دا سیٹ
وی) !!{nonenumerable}
اے؛ پر ساریاں
ٹیورنگ مشیناں
دی فہرست بن
سکدی اے، ایس
لئی ٹیورنگ
مشین نال حساب
ہو سکّن والے
فنکشناں دا
سیٹ صرف
!!{denumerable}
اے۔

اسی \emph{خاص}
فنکشن تے مسئلے
وی تعریف کر
سکدے آں جنہاں
نوں بالترتیب
ناقابلِ حساب
تے ناقابلِ فیصلہ
ثابت کیتا جا
سکدا اے۔ اک
مسئلہ اوہ اے جیہنوں
\emph{رُکّن دا مسئلہ}
کہندے نیں۔ ٹیورنگ
مشین دے حکماں
دی فہرست نال
اوس دا محدود
بیان دتا جا
سکدا اے۔ ٹیورنگ
مشین دا ایسا
بیان، یعنی
مشین دا پروگرام،
یقیناً دوجی
ٹیورنگ مشین
لئی انپٹ
بن سکدا اے۔
ایس لئی اسی
ایہناں مشیناں
بارے سوالاں
دا فیصلہ کرن
والیاں ٹیورنگ
مشیناں اُتے
غور کر سکدے
آں۔ اک خاص
دلچسپ سوال
ایہ اے:
“کی دتی گئی
ٹیورنگ مشین
انپٹ~$n$
نال شروع ہو
کے آخرکار
رُکے گی؟”
چنگا ہوندا جے
ایس سوال دا
فیصلہ کرن
والی ٹیورنگ
مشین ہوندی:
ایہنوں معیار
جانچ والی
مشین سمجھو
جیڑی دیکھدی
کہ ہور مشیناں
لامتناہی چکّر
وچ نہ پھسن۔
دلچسپ حقیقت،
جیڑی ٹیورنگ
نے ثابت کیتی،
ایہ اے کہ
ایسی مشین
موجود نہیں۔
کوئی اک ٹیورنگ
مشین ایسی
نہیں ہو سکدی
جیڑی ٹیورنگ
مشین~$M$
دے بیان تے
کسے عدد~$n$
نوں انپٹ
وچ لے کے
ہمیشہ رُکے
تے آؤٹ پٹ
$1$ یا $0$
ایس حساب نال
دیوے کہ
$M$ مشین
انپٹ $n$
نال شروع ہو
کے رُکدی
اے یا نہیں۔

جدوں ساڈے کول
خاص ناقابلِ
فیصلہ مسئلیاں
دیاں مثالاں
ہون، تاں اسی
اوہناں نال ہور
مسئلیاں دا وی
ناقابلِ فیصلہ
ہونا دکھا سکدے
آں۔ مثال دے
طور اُتے اک
مشہور ناقابلِ
فیصلہ سوال
ایہ اے: “کی
پہلے درجے دا
!!{formula}~$!A$
درست اے؟”
کوئی ٹیورنگ مشین ایہو جیہی نہیں
جیہنوں انپٹ وجوں پہلے درجے دا
!!{formula}~$!A$ دتا جاوے
تے اوہ ہر صورت رُکے تے
آؤٹ پٹ $1$
یا $0$ ایس
حساب نال
دیوے کہ $!A$
درست اے
یا نہیں۔
تاریخی طور
اُتے ایس
مسئلے نوں
مؤثر طریقے
نال حل کرن
والا طریقۂ
کار لبھݨ
دے سوال
نوں بس
“فیصلے دا
مسئلہ”
کہندے سن؛
ایس لئی
اسی کہندے
آں کہ
فیصلے دا
مسئلہ حل
نہیں ہو
سکدا۔
ٹیورنگ
تے چرچ
نے ایہ
نتیجہ لگ
بھگ اکّو
ویلے وکھ
وکھ ثابت
کیتا سی،
ایس لئی
ایہنوں
چرچ--ٹیورنگ
دا قضیہ
وی کہندے
نیں۔

\end{document}
